%!TEX root = ../dokumentation.tex
% CHEAT SHEET: every element of the template once. The PDF shows the result,
% this file the source. To remove it, delete the line %!TEX root = ../dokumentation.tex
% CHEAT SHEET: every element of the template once. The PDF shows the result,
% this file the source. To remove it, delete the line %!TEX root = ../dokumentation.tex
% CHEAT SHEET: every element of the template once. The PDF shows the result,
% this file the source. To remove it, delete the line %!TEX root = ../dokumentation.tex
% CHEAT SHEET: every element of the template once. The PDF shows the result,
% this file the source. To remove it, delete the line \input{content/beispiele}
% in dokumentation.tex.

\chapter{Examples (Cheat Sheet)}\label{ch:examples}

This chapter shows every element of the template once.
Its source is \texttt{content/beispiele.tex}.

\section{Citing}\label{sec:ex-citing}

A citation in parentheses goes at the end of the sentence \citep{shannon1948mathematical}.
With a page number it looks like this \citep[p.~42]{knuth1984texbook}.
Cited within the sentence: \citet{vaswani2017attention} introduce the Transformer architecture.
Online sources carry their URL in the bibliography \citep{overleaf2026bibtex}.
New sources go into \texttt{literatur.bib}; the bibliography lists only cited entries.

\section{Abbreviations and Glossary}\label{sec:ex-abbreviations}

The first use of \ac{API} spells out the long form; after that, only \ac{API} appears.
Glossary terms such as \gls{synctex} link to the glossary.

\section{Figures}\label{sec:ex-figures}

\autoref{fig:ex-image} shows an image; your own images go into \texttt{images/}.

\begin{figure}[htbp]
  \centering
  \includegraphics[width=0.5\textwidth]{example-image}
  \caption{Example image: give the file name without its extension}
  \label{fig:ex-image}
\end{figure}

\section{Diagrams}\label{sec:ex-diagrams}

\autoref{fig:ex-diagram} uses the TikZ style kit from \texttt{ads/header.tex}.
A single setting, \texttt{akzent}, sets the color of all diagrams.

\begin{figure}[htbp]
  \centering
  \begin{adjustbox}{max width=\textwidth}
  \begin{tikzpicture}[flow, node distance=6mm and 8mm]
    \node[fio]                   (in)   {Input};
    \node[fstage, right=of in]   (prep) {Preprocessing};
    \node[fgate,  right=of prep] (gate) {confident?};
    \node[fkey,   right=of gate] (key)  {expensive\\step};
    \node[fio,    right=of key]  (out)  {Output};
    \node[fmuted, below=of gate] (fast) {fast path};
    \draw[fedge] (in) -- (prep);
    \draw[fedge] (prep) -- (gate);
    \draw[fedge] (gate) -- node[flab, above] {no} (key);
    \draw[fcond] (gate) -- node[flab, right] {yes} (fast);
    \draw[fedge] (key) -- (out);
    \draw[fcond] (fast) -| (out);
    \node[fnote, below=1mm of in] {step 1};
  \end{tikzpicture}
  \end{adjustbox}
  \caption[Example diagram]{Example diagram with the styles \texttt{fio}, \texttt{fstage}, \texttt{fgate}, \texttt{fkey} and \texttt{fmuted}. The short title in square brackets appears in the list of figures.}
  \label{fig:ex-diagram}
\end{figure}

\section{Tables}\label{sec:ex-tables}

\autoref{tab:ex-table} uses \texttt{booktabs}: horizontal rules only, caption on top.

\begin{table}[htbp]
  \centering
  \caption{Example table}
  \label{tab:ex-table}
  \begin{tabular}{lrr}
    \toprule
    Variant & Accuracy & Time per document \\
    \midrule
    Variant A & 72.4\,\% & 3\,ms \\
    Variant B & 88.0\,\% & 40\,ms \\
    \bottomrule
  \end{tabular}
\end{table}

\section{Source Code}\label{sec:ex-code}

\autoref{lst:ex-code} shows a listing; \texttt{[language=Java]} switches the language per listing.

\begin{lstlisting}[caption={Example listing in Python}, label={lst:ex-code}]
def greet(name: str) -> str:
    # non-ASCII characters work: äöü ß
    return f"Hello, {name}!"
\end{lstlisting}

\section{Equations}\label{sec:ex-equations}

Numbered equations go into an \texttt{equation} environment; \eqref{eq:ex-f1} refers to it.
\begin{equation}
  F_1 = \frac{2 \cdot P \cdot R}{P + R}
  \label{eq:ex-f1}
\end{equation}

\section{Miscellaneous}\label{sec:ex-misc}

A footnote attaches directly to the word.\footnote{This is what a footnote looks like.}
\enquote{Quotation marks} come from \texttt{\textbackslash enquote}, matching the language.
Cross-references with \texttt{\textbackslash autoref} name the type of their target: \autoref{sec:ex-citing}, \autoref{fig:ex-image}, \autoref{tab:ex-table}.
Amounts in euros: 120\,\euro{}.

\begin{notiz}
Red draft note to yourself.
Before submitting, \texttt{\textbackslash excludecomment\{notiz\}} in \texttt{ads/header.tex} hides all notes.
\end{notiz}

% in dokumentation.tex.

\chapter{Examples (Cheat Sheet)}\label{ch:examples}

This chapter shows every element of the template once.
Its source is \texttt{content/beispiele.tex}.

\section{Citing}\label{sec:ex-citing}

A citation in parentheses goes at the end of the sentence \citep{shannon1948mathematical}.
With a page number it looks like this \citep[p.~42]{knuth1984texbook}.
Cited within the sentence: \citet{vaswani2017attention} introduce the Transformer architecture.
Online sources carry their URL in the bibliography \citep{overleaf2026bibtex}.
New sources go into \texttt{literatur.bib}; the bibliography lists only cited entries.

\section{Abbreviations and Glossary}\label{sec:ex-abbreviations}

The first use of \ac{API} spells out the long form; after that, only \ac{API} appears.
Glossary terms such as \gls{synctex} link to the glossary.

\section{Figures}\label{sec:ex-figures}

\autoref{fig:ex-image} shows an image; your own images go into \texttt{images/}.

\begin{figure}[htbp]
  \centering
  \includegraphics[width=0.5\textwidth]{example-image}
  \caption{Example image: give the file name without its extension}
  \label{fig:ex-image}
\end{figure}

\section{Diagrams}\label{sec:ex-diagrams}

\autoref{fig:ex-diagram} uses the TikZ style kit from \texttt{ads/header.tex}.
A single setting, \texttt{akzent}, sets the color of all diagrams.

\begin{figure}[htbp]
  \centering
  \begin{adjustbox}{max width=\textwidth}
  \begin{tikzpicture}[flow, node distance=6mm and 8mm]
    \node[fio]                   (in)   {Input};
    \node[fstage, right=of in]   (prep) {Preprocessing};
    \node[fgate,  right=of prep] (gate) {confident?};
    \node[fkey,   right=of gate] (key)  {expensive\\step};
    \node[fio,    right=of key]  (out)  {Output};
    \node[fmuted, below=of gate] (fast) {fast path};
    \draw[fedge] (in) -- (prep);
    \draw[fedge] (prep) -- (gate);
    \draw[fedge] (gate) -- node[flab, above] {no} (key);
    \draw[fcond] (gate) -- node[flab, right] {yes} (fast);
    \draw[fedge] (key) -- (out);
    \draw[fcond] (fast) -| (out);
    \node[fnote, below=1mm of in] {step 1};
  \end{tikzpicture}
  \end{adjustbox}
  \caption[Example diagram]{Example diagram with the styles \texttt{fio}, \texttt{fstage}, \texttt{fgate}, \texttt{fkey} and \texttt{fmuted}. The short title in square brackets appears in the list of figures.}
  \label{fig:ex-diagram}
\end{figure}

\section{Tables}\label{sec:ex-tables}

\autoref{tab:ex-table} uses \texttt{booktabs}: horizontal rules only, caption on top.

\begin{table}[htbp]
  \centering
  \caption{Example table}
  \label{tab:ex-table}
  \begin{tabular}{lrr}
    \toprule
    Variant & Accuracy & Time per document \\
    \midrule
    Variant A & 72.4\,\% & 3\,ms \\
    Variant B & 88.0\,\% & 40\,ms \\
    \bottomrule
  \end{tabular}
\end{table}

\section{Source Code}\label{sec:ex-code}

\autoref{lst:ex-code} shows a listing; \texttt{[language=Java]} switches the language per listing.

\begin{lstlisting}[caption={Example listing in Python}, label={lst:ex-code}]
def greet(name: str) -> str:
    # non-ASCII characters work: äöü ß
    return f"Hello, {name}!"
\end{lstlisting}

\section{Equations}\label{sec:ex-equations}

Numbered equations go into an \texttt{equation} environment; \eqref{eq:ex-f1} refers to it.
\begin{equation}
  F_1 = \frac{2 \cdot P \cdot R}{P + R}
  \label{eq:ex-f1}
\end{equation}

\section{Miscellaneous}\label{sec:ex-misc}

A footnote attaches directly to the word.\footnote{This is what a footnote looks like.}
\enquote{Quotation marks} come from \texttt{\textbackslash enquote}, matching the language.
Cross-references with \texttt{\textbackslash autoref} name the type of their target: \autoref{sec:ex-citing}, \autoref{fig:ex-image}, \autoref{tab:ex-table}.
Amounts in euros: 120\,\euro{}.

\begin{notiz}
Red draft note to yourself.
Before submitting, \texttt{\textbackslash excludecomment\{notiz\}} in \texttt{ads/header.tex} hides all notes.
\end{notiz}

% in dokumentation.tex.

\chapter{Examples (Cheat Sheet)}\label{ch:examples}

This chapter shows every element of the template once.
Its source is \texttt{content/beispiele.tex}.

\section{Citing}\label{sec:ex-citing}

A citation in parentheses goes at the end of the sentence \citep{shannon1948mathematical}.
With a page number it looks like this \citep[p.~42]{knuth1984texbook}.
Cited within the sentence: \citet{vaswani2017attention} introduce the Transformer architecture.
Online sources carry their URL in the bibliography \citep{overleaf2026bibtex}.
New sources go into \texttt{literatur.bib}; the bibliography lists only cited entries.

\section{Abbreviations and Glossary}\label{sec:ex-abbreviations}

The first use of \ac{API} spells out the long form; after that, only \ac{API} appears.
Glossary terms such as \gls{synctex} link to the glossary.

\section{Figures}\label{sec:ex-figures}

\autoref{fig:ex-image} shows an image; your own images go into \texttt{images/}.

\begin{figure}[htbp]
  \centering
  \includegraphics[width=0.5\textwidth]{example-image}
  \caption{Example image: give the file name without its extension}
  \label{fig:ex-image}
\end{figure}

\section{Diagrams}\label{sec:ex-diagrams}

\autoref{fig:ex-diagram} uses the TikZ style kit from \texttt{ads/header.tex}.
A single setting, \texttt{akzent}, sets the color of all diagrams.

\begin{figure}[htbp]
  \centering
  \begin{adjustbox}{max width=\textwidth}
  \begin{tikzpicture}[flow, node distance=6mm and 8mm]
    \node[fio]                   (in)   {Input};
    \node[fstage, right=of in]   (prep) {Preprocessing};
    \node[fgate,  right=of prep] (gate) {confident?};
    \node[fkey,   right=of gate] (key)  {expensive\\step};
    \node[fio,    right=of key]  (out)  {Output};
    \node[fmuted, below=of gate] (fast) {fast path};
    \draw[fedge] (in) -- (prep);
    \draw[fedge] (prep) -- (gate);
    \draw[fedge] (gate) -- node[flab, above] {no} (key);
    \draw[fcond] (gate) -- node[flab, right] {yes} (fast);
    \draw[fedge] (key) -- (out);
    \draw[fcond] (fast) -| (out);
    \node[fnote, below=1mm of in] {step 1};
  \end{tikzpicture}
  \end{adjustbox}
  \caption[Example diagram]{Example diagram with the styles \texttt{fio}, \texttt{fstage}, \texttt{fgate}, \texttt{fkey} and \texttt{fmuted}. The short title in square brackets appears in the list of figures.}
  \label{fig:ex-diagram}
\end{figure}

\section{Tables}\label{sec:ex-tables}

\autoref{tab:ex-table} uses \texttt{booktabs}: horizontal rules only, caption on top.

\begin{table}[htbp]
  \centering
  \caption{Example table}
  \label{tab:ex-table}
  \begin{tabular}{lrr}
    \toprule
    Variant & Accuracy & Time per document \\
    \midrule
    Variant A & 72.4\,\% & 3\,ms \\
    Variant B & 88.0\,\% & 40\,ms \\
    \bottomrule
  \end{tabular}
\end{table}

\section{Source Code}\label{sec:ex-code}

\autoref{lst:ex-code} shows a listing; \texttt{[language=Java]} switches the language per listing.

\begin{lstlisting}[caption={Example listing in Python}, label={lst:ex-code}]
def greet(name: str) -> str:
    # non-ASCII characters work: äöü ß
    return f"Hello, {name}!"
\end{lstlisting}

\section{Equations}\label{sec:ex-equations}

Numbered equations go into an \texttt{equation} environment; \eqref{eq:ex-f1} refers to it.
\begin{equation}
  F_1 = \frac{2 \cdot P \cdot R}{P + R}
  \label{eq:ex-f1}
\end{equation}

\section{Miscellaneous}\label{sec:ex-misc}

A footnote attaches directly to the word.\footnote{This is what a footnote looks like.}
\enquote{Quotation marks} come from \texttt{\textbackslash enquote}, matching the language.
Cross-references with \texttt{\textbackslash autoref} name the type of their target: \autoref{sec:ex-citing}, \autoref{fig:ex-image}, \autoref{tab:ex-table}.
Amounts in euros: 120\,\euro{}.

\begin{notiz}
Red draft note to yourself.
Before submitting, \texttt{\textbackslash excludecomment\{notiz\}} in \texttt{ads/header.tex} hides all notes.
\end{notiz}

% in dokumentation.tex.

\chapter{Examples (Cheat Sheet)}\label{ch:examples}

This chapter shows every element of the template once.
Its source is \texttt{content/beispiele.tex}.

\section{Citing}\label{sec:ex-citing}

A citation in parentheses goes at the end of the sentence \citep{shannon1948mathematical}.
With a page number it looks like this \citep[p.~42]{knuth1984texbook}.
Cited within the sentence: \citet{vaswani2017attention} introduce the Transformer architecture.
Online sources carry their URL in the bibliography \citep{overleaf2026bibtex}.
New sources go into \texttt{literatur.bib}; the bibliography lists only cited entries.

\section{Abbreviations and Glossary}\label{sec:ex-abbreviations}

The first use of \ac{API} spells out the long form; after that, only \ac{API} appears.
Glossary terms such as \gls{synctex} link to the glossary.

\section{Figures}\label{sec:ex-figures}

\autoref{fig:ex-image} shows an image; your own images go into \texttt{images/}.

\begin{figure}[htbp]
  \centering
  \includegraphics[width=0.5\textwidth]{example-image}
  \caption{Example image: give the file name without its extension}
  \label{fig:ex-image}
\end{figure}

\section{Diagrams}\label{sec:ex-diagrams}

\autoref{fig:ex-diagram} uses the TikZ style kit from \texttt{ads/header.tex}.
A single setting, \texttt{akzent}, sets the color of all diagrams.

\begin{figure}[htbp]
  \centering
  \begin{adjustbox}{max width=\textwidth}
  \begin{tikzpicture}[flow, node distance=6mm and 8mm]
    \node[fio]                   (in)   {Input};
    \node[fstage, right=of in]   (prep) {Preprocessing};
    \node[fgate,  right=of prep] (gate) {confident?};
    \node[fkey,   right=of gate] (key)  {expensive\\step};
    \node[fio,    right=of key]  (out)  {Output};
    \node[fmuted, below=of gate] (fast) {fast path};
    \draw[fedge] (in) -- (prep);
    \draw[fedge] (prep) -- (gate);
    \draw[fedge] (gate) -- node[flab, above] {no} (key);
    \draw[fcond] (gate) -- node[flab, right] {yes} (fast);
    \draw[fedge] (key) -- (out);
    \draw[fcond] (fast) -| (out);
    \node[fnote, below=1mm of in] {step 1};
  \end{tikzpicture}
  \end{adjustbox}
  \caption[Example diagram]{Example diagram with the styles \texttt{fio}, \texttt{fstage}, \texttt{fgate}, \texttt{fkey} and \texttt{fmuted}. The short title in square brackets appears in the list of figures.}
  \label{fig:ex-diagram}
\end{figure}

\section{Tables}\label{sec:ex-tables}

\autoref{tab:ex-table} uses \texttt{booktabs}: horizontal rules only, caption on top.

\begin{table}[htbp]
  \centering
  \caption{Example table}
  \label{tab:ex-table}
  \begin{tabular}{lrr}
    \toprule
    Variant & Accuracy & Time per document \\
    \midrule
    Variant A & 72.4\,\% & 3\,ms \\
    Variant B & 88.0\,\% & 40\,ms \\
    \bottomrule
  \end{tabular}
\end{table}

\section{Source Code}\label{sec:ex-code}

\autoref{lst:ex-code} shows a listing; \texttt{[language=Java]} switches the language per listing.

\begin{lstlisting}[caption={Example listing in Python}, label={lst:ex-code}]
def greet(name: str) -> str:
    # non-ASCII characters work: äöü ß
    return f"Hello, {name}!"
\end{lstlisting}

\section{Equations}\label{sec:ex-equations}

Numbered equations go into an \texttt{equation} environment; \eqref{eq:ex-f1} refers to it.
\begin{equation}
  F_1 = \frac{2 \cdot P \cdot R}{P + R}
  \label{eq:ex-f1}
\end{equation}

\section{Miscellaneous}\label{sec:ex-misc}

A footnote attaches directly to the word.\footnote{This is what a footnote looks like.}
\enquote{Quotation marks} come from \texttt{\textbackslash enquote}, matching the language.
Cross-references with \texttt{\textbackslash autoref} name the type of their target: \autoref{sec:ex-citing}, \autoref{fig:ex-image}, \autoref{tab:ex-table}.
Amounts in euros: 120\,\euro{}.

\begin{notiz}
Red draft note to yourself.
Before submitting, \texttt{\textbackslash excludecomment\{notiz\}} in \texttt{ads/header.tex} hides all notes.
\end{notiz}
